\begin{statement}{A.4}{Proved}{All first-cover models, Jacobians, and cusp divisors}
Put
\[
 (a,b)\in\mathcal T=\{(0,-1),(6,1),(-6,1)\},\qquad
 E_{a,b}:y^2=x^3+ax^2+bx.
\]
Each $E_{a,b}$ has a degree-two isogeny to $E$, with kernel
$\{O,(0,0)\}$. The three geometric $S_3$-covers and three geometric
$D_8$-covers of $U$ have smooth compactifications
\[
 C_{3,a,b}:v^2=w^7+aw^4+bw,\qquad
 C_{4,a,b}:v^2=w^8+aw^4+b,                         \tag{A.8}
\]
respectively. Their maps to $E_{a,b}$ are $(w^3,wv)$ and
$(w^4,w^2v)$; compose with the indicated isogeny to $E$.
All six curves have genus three. The $C_3$ cusp set is
$\{P_0,P_\infty\}$, and the $C_4$ cusp set consists of the two
points over $w=0$ and the two over $w=\infty$.

For the $C_3$ models take $k=-1$ if $(a,b)=(0,-1)$ and $k=1$
otherwise, so $k^3=b$. For the $C_4$ models choose $c^4=b$.
The following are geometric isogeny decompositions and quotient models:
\[
\begin{array}{c|c|c}
\text{curve}&\text{elliptic complementary factor}&\text{Jacobian}\\ \hline
C_{3,a,b}&F_{3,a,k}:u^2=t^3-3kt+a&E_{a,b}\times F_{3,a,k}^{\,2}\rule{0pt}{13pt}\\
C_{4,a,b}&F_{4,a,c}:u^2=t^4-4ct^2+a+2c^2&E_{a,b}\times F_{4,a,c}^{\,2}
\end{array}                                                     \tag{A.9}
\]
The genus-two quotients by a reflection have equations
\[
\begin{split}
 D_{3,a,k}:&\quad z^2=(t^2-4k)(t^3-3kt+a),\\
 D_{4,a,c}:&\quad z^2=(t^2-4c)(t^4-4ct^2+a+2c^2).
\end{split}                                                     \tag{A.10}
\]
Their Jacobians are isogenous to $E_{a,b}\times F_{3,a,k}$ and
$E_{a,b}\times F_{4,a,c}$, respectively. The $D_3$ quotient has
one cusp at infinity; the $D_4$ quotient has two. For $D_8$ the choices
$c$ and $ic$ give the two reflection conjugacy classes; $c$ and $-c$
give isomorphic quotients.

The $E$-isotypic dimensions in (A.9) are three for $C_{3,0,-1}$
and one for each of the other five curves. In (A.10) they are two
for $D_{3,0,-1}$ and one otherwise. In more explicit terms,
\[
\begin{array}{c|c}
\text{complementary factor}&j\\ \hline
F_{3,0,-1}&1728\rule{0pt}{13pt}\\
F_{3,\pm6,1}&-216\\
F_{4,0,c}&8000\\
F_{4,\pm6,c}&10976\text{ or }128
\end{array}                                                     \tag{A.11}
\]
The last two values are isogenous complementary factors, not additional
$E$ factors.

The unique geometric $Q_8$-cover has compactification
\[
 C_Q:X^4+Y^4=2Z^4,
\quad x_E=(xy)^2,\quad y_E=\frac{i}{2}xy(x^4-y^4),
\quad x=X/Z,\ y=Y/Z.                                            \tag{A.12}
\]
Its four cusps are $Z=0$, and $\operatorname{Jac}(C_Q)\sim E^3$.
It has no genus-two quotient, even by an automorphism outside its
$Q_8$ deck group. Thus an expectation of a genus-two quotient for
this particular cover must be corrected.

The groups generated by degree-zero cusp divisor classes are exactly
\[
\begin{array}{c|c}
 C_{3,a,b}&\Z/2\rule{0pt}{13pt}\\
 D_{3,a,k}&0\\
 C_{4,a,b}&(\Z/4)^2\\
 D_{4,a,c}&\Z/4\\
 C_Q&(\Z/4)^2.
\end{array}                                                     \tag{A.13}
\]
In particular all projected cusp divisors in A.3 are torsion for
these covers. The dimensions of their potentially $\Delta$-fixed
rational homology are, in the same order, $7$ for $C_{3,0,-1}$ and
$3$ for the other $C_3$ curves; $4$ for $D_{3,0,-1}$ and $2$ for
the other $D_3$ curves; $5$ for every $C_4$; $3$ for every $D_4$;
and $9$ for $C_Q$. Their potentially Galois-stable rational lines
are then exactly the boundary and CM lines described in A.3.
\end{statement}
\begin{proof}
We give constructions and the divisor calculations, so no computer-algebra
decomposition of a Jacobian is assumed. The standard degree-two isogeny
with kernel $(0,0)$ is
\[
 (x,y)\longmapsto
 \left(x+a+\frac b x,\ y\left(1-\frac b{x^2}\right)\right)
\]
from $E_{a,b}$ to $y'^2=x'^3-2ax'^2+(a^2-4b)x'$. For
$a=\pm6,b=1$, the substitution $x' =4x_E+2a/3$, $y'=8y_E$
identifies the target with $E$. For $a=0,b=-1$, use
$x_E=-ix'/2$, $y_E=-(1+i)y'/4$; this is the isogeny $[1+i]$
with the indicated coordinates. Direct substitution proves both
covering maps in (A.8).

Let $\omega^3=1$, $\omega\ne1$. The $C_3$ deck generators are
\[
 \tau(w,v)=(\omega w,\omega^{-1}v),\qquad
 j(w,v)=(k/w,-k^2v/w^4).
\]
They satisfy $\tau^3=j^2=1$, $j\tau j=\tau^{-1}$.
The second map induces translation by $(0,0)$ on $E_{a,b}$.
For $C_4$, the corresponding generators are
\[
 \tau(w,v)=(iw,-v),\qquad j(w,v)=(c/w,-c^2v/w^4),
\]
and satisfy $\tau^4=j^2=1$, $j\tau j=\tau^{-1}$.
The polynomials in (A.8) have simple roots for our three parameter
pairs, so their indicated smooth projective models have genus three.
The only points above the removed kernel of the isogeny are exactly
the displayed cusps. The maps have total degrees six and eight, and
the displayed deck groups have those orders, hence these are the
required Galois covers.

For completeness they exhaust the covers, rather than merely giving
examples. Surjections from the free rank-two geometric group to the
three finite groups are generating pairs modulo group automorphisms.
The executed independent program \texttt{a\_first\_covers.py}
enumerates $18,24,24$ generating pairs in $S_3,D_8,Q_8$.
Their automorphism groups have orders $6,8,24$, and the actions on
generating pairs are free. Thus the numbers of covers are $3,3,1$.
In each of the first two cases the three displayed covers have the
three different degree-two intermediate covers of $E\setminus O$,
corresponding to the three nonzero two-torsion subgroups, so they are
distinct. Although the two sources with $a=6$ and $a=-6$ have equal
$j$, their dual kernels on $E$ are the two different nonzero points
other than $(0,0)$; the source isomorphism covers $[i]$ on $E$, not
the identity. This also identifies the unique $[i]$-stable $S_3$ cover.

Let $h(w,v)=(w,-v)$ be the hyperelliptic involution. On $C_3$ set
\[
 t=w+k/w,\qquad u=v/w^2,\qquad z=(w-k/w)u.
\]
The map $(t,u)$ is the quotient by $hj$, and $(t,z)$ the quotient
by $j$. The identities
\[
 w^3+k^3/w^3=t^3-3kt,\qquad (w-k/w)^2=t^2-4k
\]
give the first equations in (A.9)--(A.10). Similarly, on $C_4$ use
$t=w+c/w$, $u=v/w^2$, and $z=(w-c/w)u$. The identity
\[
 w^4+c^4/w^4=t^4-4ct^2+2c^2
\]
gives the second equations. Each $j$ is fixed-point-free: a fixed
point would have $w^2=k$ or $w^2=c$ and $v=0$, which our parameter
values exclude. Riemann--Hurwitz therefore gives genus two for its
quotient. The quotient by $hj$ has genus one.

Here is also a direct explanation of the isogeny factors. The quotient
by the rotation group has Jacobian $E_{a,b}$. Its differential is the
one-dimensional rotation-invariant summand of
$\langle dw/v,w\,dw/v,w^2dw/v\rangle$. The remaining two-dimensional
summand is the standard rational representation of $S_3$, respectively
$D_8$. Its simple rational group-algebra component is $M_2(\Q)$.
The two reflection projectors give isogenous one-dimensional abelian
factors; the quotient by $hj$ supplies one of them. Clearing the
denominators of the rational idempotents gives the isogenies in (A.9).
The $j$-invariant part similarly gives (A.10). Equivalently the images
of the elliptic quotient maps span the differential spaces, which
makes the induced map of abelian varieties an isogeny.

The $j$ values follow either from the cubic equation or from
\[
 j(u^2=t^4+Bt^2+C)=
       \frac{16(B^2+12C)^3}{C(B^2-4C)^2}.
\]
After scaling $t$ by a square root of $c$, the $C_4$ complements
have $B=-4$ and $C=2+a/c^2$, hence $C=2,8,-4$.
It remains to justify that none of the complements with $j\ne1728$
is geometrically isogenous to $E$; inequality of $j$ alone would
not justify that assertion. For $F_{3,6,1}$ the trace at $5$ is
$-1$, with Frobenius field $\Q(\sqrt{-19})$, whereas $E$ has
trace $-2$ and field $\Q(i)$. The eigenvalue quotients are not
roots of unity, so all positive powers still generate those
different quadratic fields. An isogeny over any finite extension
would identify some powers of these Frobenius polynomials, which
is impossible. The $a=-6$ factor is a quadratic twist and gives
the same conclusion. For $F_{4,0,c}$, use its normalized quartic
at $3$: its trace is $2$, hence its reduction is ordinary, while
$E$ has supersingular reduction. For the other two quartics use
$7$, where their traces are $-4$ and $4$, while $E$ has trace zero
and is supersingular. All these primes are primes of good reduction
for the indicated equations. Geometric isogeny preserves the
ordinary/supersingular distinction. The independent point counts
are recorded in \texttt{a\_first\_covers.json}.

For (A.12) the automorphisms
\[
 \alpha(x,y)=(ix,-iy),\qquad \beta(x,y)=(y,-x)
\]
satisfy $\alpha^2=\beta^2$, $\alpha^4=1$ and
$\beta\alpha\beta^{-1}=\alpha^{-1}$, and generate $Q_8$.
Their common central involution has exactly the four points $Z=0$
as fixed points; the other nonidentity elements outside the center
have none. The invariant functions in (A.12) satisfy $y_E^2=x_E^3-x_E$
and have degree eight, so give the quotient. The three sign
involutions on $X,Y,Z$ yield elliptic quotients with $j=1728$.
Their differential projectors span the three-dimensional space,
giving $\operatorname{Jac}(C_Q)\sim E^3$.

The smooth plane quartic $C_Q$ is nonhyperelliptic, and its canonical
embedding is this plane model. Every involution extends to a
projective linear involution of the plane and fixes a line. That
line meets the curve in four distinct points: tangency would give
derivative one for the involution at a fixed point of the curve,
forcing the involution to be the identity in characteristic zero.
Thus no involution is free. A map from genus three to genus two
would, by Riemann--Hurwitz, have degree two and be unramified.
Such a map cannot exist. This proves the assertion about genus-two
quotients without relying merely on the list of subgroups of $Q_8$.

We finish with the exact cusp-divisor groups. On $C_3$,
$\operatorname{div}(w)=2(P_0-P_\infty)$. The class is nonzero:
a function with divisor $P_0-P_\infty$ would give a degree-one map
from a positive-genus curve to $\mathbf P^1$. The $D_3$ quotient
has only one cusp.

For $C_4$ write $s=c^2$, so $s^2=b$, and label the cusps
$A=(0,s)$, $B=(0,-s)$, $C=(\infty,+)$, $D=(\infty,-)$.
The order-four automorphism $(w,v)\mapsto(iw,v)$ fixes all four.
Its quotient is the rational conic $v^2=q^2+aq+b$, $q=w^4$,
with parameter
\[
 T=(v-s)/w^4,\qquad w^4=\frac{a-2sT}{T^2-1}.
\]
The four cusp images are $a/(2s),\infty,1,-1$; they are distinct.
Pullbacks of rational functions on this line kill four times each
cusp difference, and
$\operatorname{div}(w)=A+B-C-D$ is an additional relation.
These are all relations. Indeed a function whose divisor is
supported on the four cusps is an eigenfunction for the order-four
automorphism, since that automorphism fixes its divisor; its
quotient with its transform is constant. It therefore has the form
$w^e R(T)$, $0\le e<4$, and its valuations are exactly those
generated by the stated relations. The resulting Smith form is
$(1,4,4)$ on the degree-zero divisor lattice, so the group is
$(\Z/4)^2$.

For $D_4$ write $Q_+,Q_-$ for the two points at infinity and
$\delta=[Q_+-Q_-]$. Pushforward from $C_4$ shows $4\delta=0$.
Let $W_+,W_-$ be its Weierstrass points $t^2=4c$.
The connected etale double cover $C_4\to D_4$ corresponds to
the nonzero two-torsion class
\[
 \eta=[W_++W_--Q_+-Q_-],
\]
as is seen by adjoining $\sqrt{t^2-4c}$. The identity
\[
 z^2-(t^3-4ct)^2=(a+2c^2)(t^2-4c)
\]
shows that
$\operatorname{div}(z-t^3+4ct)=W_++W_-+Q_+-3Q_-$,
with the labels chosen by $z/t^3=+1$ at $Q_+$.
Consequently $\eta=-2\delta\ne0$, so $\delta$ has exact order four.

Finally the four points $Z=0$ on $C_Q$ are the totally ramified
points of the cyclic fourth-root cover
\[
 (Z/Y)^4=((X/Y)^4+1)/2.
\]
The same eigenfunction argument gives all relations in their
degree-zero divisor group: four times each difference, and the
relation whose coefficient vector in the basis $P_j-P_4$ is
$(1,1,1)$. The latter is also visible from the line $Z=0$ and a
hyperflex tangent line, whose divisor is $4P_4$. Its Smith form
is again $(1,4,4)$, proving (A.13).

All divisor assertions used in A.3 are therefore torsion assertions,
not inferred just from the CM decompositions. Applying A.3 now
gives each listed dimension as $2\dim J_E+|S|-1$.
\end{proof}

\begin{statement}{A.10}{Proved}{The complete first-cover calculation}
For every first cover and genus-two quotient specified in Task A.4,
at every prime and for every stated packet, the subgroup fixed by
the permitted open subgroup of $\Delta$ is exactly the selected
cusp inertia. Every nontrivial closed procyclic subgroup with
open arithmetic normalizer is contained in a cusp inertia group,
up to conjugacy. Together with A.4 this completes Task A.4.

The explicit models, isogenies, cusp classes, and rational homology
of Task A.4 are determined by A.4. A.22 determines the full fixed subgroups for covers with
$\ell$-group Galois closure, including every $D_8/Q_8$ case at
two. A.27 determines the stable genus-two symmetric-cover fixed groups
for every original packet and every odd prime. A.30 determines the
stable genus-three symmetric-cover fixed groups for every original
packet at every odd prime, including three. A.31 supplies the
two-adic comparison and settles both stable symmetric curves at
two as well. A.32 determines the two nonstable symmetric
genus-two quotient families at every prime.
A.33 settles the two nonstable symmetric genus-three covers at two
and reduces their odd-prime leading fixed classes to
$\mathbb L(t,\Omega)$ in every degree.
A.35 removes the entire proper image in the quaternion case at every odd
prime, and A.36 determines its full punctured fixed group at every prime.
A.37 determines all three dihedral genus-three fixed groups at every prime.
A.38 determines the two nonstable symmetric genus-three fixed groups
at every prime, using the weighted cyclotomic criterion.
A.39 determines every dihedral genus-two reflection quotient at
every prime, including both reflection classes of each cover.
A.14 and A.25 settle the procyclic normalizer classification for every
displayed genus-three cover and every displayed genus-two quotient,
at all primes. The corresponding universal assertion in Task A.5
remains Open. Potentially fixed homology is not asserted to lift to
fixed elements of $H^{(\ell)}$.
\end{statement}
\begin{proof}[Proof and comparison with the earlier reductions]
The displayed cases exhaust the covers and quotients in A.4.
The cited theorems give genuine actions with the chosen cusp lift,
and include passage to every permitted open subgroup of $\Delta$.
Their fixed-subgroup equalities, together with A.14 and A.25 for
the normalizer assertion, prove the statement. We record why the
abelian calculations alone had not sufficed.

For the $[i]$-stable $S_3$ cover the first geometric claims in Lead 3
are proved: it is $v^2=w^7-w$, its cusp class has order two, and
its Jacobian is isogenous to $E^3$. Thus its seven-dimensional
rational open-curve homology becomes fixed by an open subgroup of
$\Delta$. This is a statement about the associated abelian quotient.
The hyperelliptic involution acts as $-1$ on compact-curve homology
but as $+1$ on the boundary generator, so a parity argument treating
all seven generators alike would be invalid.

A.28 computes the actual stable marked graph at two: three elliptic
tails attached to a rational component containing both cusps.
Its separating monodromy and a CM Frobenius conjugate reduce every
possible fixed leading class to $\mathbb L(t,\Omega)$. A.29 proves
that any nonzero first Johnson matrix kills every class in that
algebra except $t$, without an invertibility assumption. A.30 computes the actual primitive cusp-path extension using the
universal Collino cycle and its elliptic norm. Its nonzero regulator
is detected by the unit $2+\sqrt3$, so the possible zero-Johnson
packet does not exist. A.31 replaces tame monodromy at two by
Fontaine's monodromy on the crystalline realization of the global
orbit span of a fixed class. Its explicit logarithmic tangential
comparison handles arbitrary packets over ramified local fields.
The stable symmetric-cover fixed groups are therefore determined
at every prime.

For the other first covers, the non-$E$ factors in (A.11) already
exclude their proper homology from the potentially fixed subspace,
but do not identify the nonabelian kernel. A.1 and A.9 completely
settle which fixed quotient elements admit a fixed full-profinite
lift: exactly those in the image of the specified cusp inertia.
They do not decide whether additional quotient-fixed elements exist.
The last symmetric case was a nonstable $S_3$
cover at an odd prime. A.33 gives its three-tail graph and
uses the derived special linear action on the non-CM factors
to reduce the kernel to $\mathbb L(t,\Omega)$. The primitive
first Johnson class takes values in a symmetric-square module;
the CM elliptic norm in A.30 does not evaluate it.
Instead, A.38 uses the full Galois action on the actual fixed
subgroup: the Tate-graded weighted kernel forces the cyclotomic
kernel to act trivially on that subgroup, after which proper-cover
homology and Ribet's theorem exclude every nonperipheral element.
A.39 applies the same criterion to the remaining dihedral
genus-two quotients. Thus no unevaluated first Johnson class
remains in Task A.4. The assertion on arbitrary finite covers,
where these explicit weighted kernels are unavailable, remains
the universal open problem of Task A.5.
\end{proof}
